\section{Why Check-Then-Update Is Unsafe}
A sequence that first reads \code{available\_quantity} and later updates it can oversell because concurrent transactions may observe the same pre-update value.

\section{Atomic Conditional Inventory Update}
\begin{lstlisting}[language=SQL,caption={Atomic inventory reservation}]
UPDATE ticket_inventories
SET available_quantity = available_quantity - :quantity
WHERE ticket_category_id = :ticket_category_id
  AND available_quantity >= :quantity
RETURNING available_quantity;
\end{lstlisting}

A returned row means the reservation succeeded. No returned row means there was insufficient inventory. The correctness decision is made by PostgreSQL rather than by a stale application read.

\section{Concurrency Strategy}
\begin{table}[H]
\centering
\caption{Concurrency hazards and controls}
\begin{tabularx}{\textwidth}{p{0.28\textwidth}Y}
\toprule
\textbf{Hazard} & \textbf{Control} \\
\midrule
Two customers compete for final ticket & Atomic conditional update \\
Booking contains multiple categories & Deterministic category order + one DB transaction \\
Voucher and inventory have different outcomes & Same transaction; any failure rolls back both \\
Admin/worker changes stale booking state & Conditional status update based on current state \\
Multiple expiry workers & Row claim with \code{FOR UPDATE SKIP LOCKED} (optional plus point) \\
\bottomrule
\end{tabularx}
\end{table}

\section{Expected Correctness Evidence}
A representative concurrency test seeds 10 available tickets and launches 100 simultaneous booking attempts. The expected result is exactly 10 successful bookings, 90 controlled sold-out responses, and final inventory of zero---never a negative value.
